\documentclass[a4paper]{article}

\usepackage[fontset=fandol]{ctex}
\usepackage{amsmath, amssymb}
\usepackage{graphicx}
\usepackage[margin=2.5cm]{geometry}
\usepackage[most]{tcolorbox}
\usepackage{etoolbox}
\usepackage{listings}
\usepackage{hyperref}
\usepackage{booktabs}
\usepackage{subcaption}
\usepackage{float}
\usepackage{tikz}
\IfFileExists{pgfplots.sty}{
    \usepackage{pgfplots}
    \pgfplotsset{compat=1.18}
}{}

\newtcolorbox{knowledgebox}[1]{
    enhanced, colback=blue!5!white, colframe=blue!75!black, colbacktitle=blue!75!black,
    coltitle=white, fonttitle=\bfseries, title=#1, attach boxed title to top left={yshift=-2mm, xshift=2mm},
    boxrule=1pt, sharp corners
}

\newtcolorbox{importantbox}[1]{
    enhanced, colback=yellow!10!white, colframe=yellow!80!black, colbacktitle=yellow!80!black,
    coltitle=black, fonttitle=\bfseries, title=#1, sharp corners
}

\newtcolorbox{warningbox}[1]{
    enhanced, colback=red!5!white, colframe=red!75!black, colbacktitle=red!75!black,
    coltitle=white, fonttitle=\bfseries, title=#1, sharp corners
}

\lstset{
    language=Python,
    basicstyle=\ttfamily\small,
    keywordstyle=\color{blue},
    stringstyle=\color{red!60!black},
    commentstyle=\color{green!60!black},
    breaklines=true,
    frame=single,
    numbers=left,
    numberstyle=\tiny\color{gray},
    captionpos=b,
    extendedchars=false
}

\newcommand{\notetitle}{CS224R 课程笔记：Lecture 1 Class Intro}
\newcommand{\noteauthors}{Codex}
\newcommand{\notedate}{2026-04-13}
\newcommand{\videochannel}{CS224R / Stanford}
\newcommand{\videopublishdate}{课程讲次日期：2025-04-02}
\newcommand{\videoduration}{未在本地元数据中提供}
\newcommand{\videourl}{https://www.youtube.com/watch?v=EvHRQhMX7_w}
\newcommand{\videocoverpath}{cover.jpg}

\begin{document}

\begin{titlepage}
\centering
{\Large 课程笔记\par}
\vspace{1.2cm}
{\huge\bfseries \notetitle\par}
\vspace{0.8cm}
{\large \noteauthors\par}
\vspace{0.3cm}
{\large \notedate\par}
\vspace{1.2cm}

\ifdefempty{\videocoverpath}{
{\small 请在生成笔记时填入视频封面图的本地路径。\par}
}{
\includegraphics[width=0.82\textwidth,height=0.45\textheight,keepaspectratio]{\videocoverpath}\par
}

\vfill
\begin{tcolorbox}[width=0.9\textwidth, colback=black!2!white, colframe=black!60, sharp corners]
\textbf{视频作者/频道}：\videochannel\par
\textbf{发布日期}：\videopublishdate\par
\textbf{视频时长}：\videoduration\par
\textbf{视频链接}：\href{\videourl}{\nolinkurl{\videourl}}
\end{tcolorbox}
\end{titlepage}

\tableofcontents
\newpage

\section{课程定位与本讲学习目标}

这节课一半是课程导论，一半是真正的技术开篇。讲者先把课程的行政与学习安排讲清楚，再用这段时间建立整门课的共同语言：什么是深度强化学习，为什么它和普通机器学习不一样，课程后续会围绕哪些算法路线展开。

\subsection{课程想让学生获得什么能力}

从课件和口头说明看，这门课的核心目标不是死记某几个算法，而是形成两层能力：
\begin{itemize}
\item 第一层是概念能力：理解深度强化学习到底在解决什么问题，常见对象包括策略、奖励、轨迹、状态与观测。
\item 第二层是实现能力：不仅要会读论文，还要能把已有方法真正实现出来，并理解不同方法的使用前提与局限。
\end{itemize}

因此课程不会试图覆盖强化学习所有理论分支，而是聚焦于最常见、最能落到深度神经网络上的方法族。讲者明确提到 imitation learning、offline / online RL、model-free / model-based RL、multi-task / meta RL，以及 RL for LLMs 和 RL for robots，这相当于给出了整门课的主题地图。

\subsection{课程组织、资源与作业安排}

本讲的 logistics 不是纯粹行政信息，因为它直接决定后续学习节奏。课程页面、Canvas、Ed、Gradescope 与 office hours 是日常协作基础；作业和项目则对应两条能力培养线：
\begin{itemize}
\item 四次两周制作业，占总评一半，重点是用 PyTorch 实作算法并在物理模拟、导航等环境中跑实验。
\item 期末项目占另一半，可以做自定义项目，也可以做默认项目。默认项目在这一年新增，主题与用 RL 微调 LLM 有关。
\end{itemize}

讲者还强调了先修要求：机器学习、基础深度学习、强化学习入门都默认掌握，原因不是要设置门槛，而是这门课会快速进入算法和系统实现，不会停留在最基础的概率与梯度下降复习上。

\begin{knowledgebox}{为什么第一讲仍然要花时间讲课程结构}
因为深度强化学习的学习曲线很陡。作业、项目、office hours 和 PyTorch tutorial 这些安排，本质上是在为后面高强度的实现型内容做缓冲。
\end{knowledgebox}

\subsection{本讲的技术学习目标}

讲者把本讲目标总结得很具体：
\begin{itemize}
\item 什么是 deep reinforcement learning；
\item 怎样表示行为，使其可以被优化；
\item 怎样把一个真实问题写成强化学习问题。
\end{itemize}

也就是说，今天并不是讲某个具体算法，而是在搭建最基础的建模框架。只要这一步不牢，后面无论是 imitation learning、policy gradients 还是 actor-critic，都会变成符号堆积。

\subsection{本章小结}

第一部分的重点是校准预期：CS224R 讲的是可以落到深度网络上的强化学习方法，目标是让学生具备理解与实现现有乃至新兴方法的能力。课程 logistics 在这里并非边角料，而是为了支撑后续的作业节奏与项目实践。

\section{为什么要学深度强化学习}

\subsection{超越监督学习：预测会改变未来}

讲者用一个非常直接的对比来切入：监督学习通常是在固定数据集上学习 \(f(x)\approx y\)，而强化学习要学的是行为 \(\pi(a\mid s)\)。这个差别看似只是符号不同，实质却是问题结构的不同。

\begin{figure}[H]
\centering
\includegraphics[width=\textwidth]{pdf_pages/page-08.png}
\caption{监督学习与强化学习的根本差别：一个学习静态输入到标签的映射，一个学习会改变未来观测的数据生成行为。}
\end{figure}

在监督学习里，模型通常被直接告知正确答案，而且样本经常可以被近似看作独立同分布。强化学习里则不同：
\begin{itemize}
\item 反馈往往是间接的，不会直接告诉系统“当前应该输出什么动作”；
\item 动作会改变下一步看到的状态或观测，因此数据分布会随着策略变化而变化；
\item 我们优化的不再只是单步预测精度，而是长期后果。
\end{itemize}

讲者用一句话概括得很准：AI model predictions have consequences。也正因为预测有后果，才需要一种把“行为”和“长期结果”绑定起来的学习框架。

\subsection{应用面广：从机器人到 LLM}

讲者随后列出了一组典型场景，说明强化学习并不是只服务于机器人或游戏。机器人控制、聊天系统、游戏智能体、自动驾驶、网页代理都可以看成“顺序决策”问题：系统不断接收信息并做动作，动作又会影响接下来可见的信息。

\begin{figure}[H]
\centering
\includegraphics[width=0.95\textwidth]{pdf_pages/page-17.png}
\caption{课程把“为什么学习深度强化学习”首先归结为超越静态 \((x,y)\) 学习范式。}
\end{figure}

这部分还给出了若干更具体的工业与研究实例：腿足机器人、机械臂操作、复杂博弈、交通控制、图像生成模型的对齐，以及现代大语言模型的后训练。其共同点在于，系统不是只要把一个标签猜对，而是要持续地产生高质量行为。

\begin{importantbox}{本课程的一个重要立场}
强化学习在现代 AI 里并不是边缘课题。尤其当系统需要和环境交互、需要从反馈中持续改进，或者需要围绕复杂目标进行优化时，RL 会自然出现。
\end{importantbox}

\subsection{从智能视角理解 RL}

讲者给出的第三个理由更加基础：learning from experience seems fundamental to intelligence。一个系统如果只能照着已有标注做预测，而不能在试错中改进，就很难具备真正开放场景中的适应性。

从这个角度看，强化学习之所以重要，不仅因为它在工业上有用，还因为它提供了一个形式化框架，去描述“练习后变好”“从结果中反思”“在不确定环境里试错”这些更接近智能本质的过程。

\subsection{开放研究问题仍然很多}

第一讲没有回避强化学习的困难。讲者明确指出，即便已经有很多成功应用，深度强化学习依然充满挑战，例如：
\begin{itemize}
\item 如何定义或学习合适的 reward；
\item 如何让智能体泛化到不同场景；
\item 如何在数据昂贵、探索危险、反馈稀疏时稳定学习；
\item 如何减少“幕后大量人工支持”的成本。
\end{itemize}

这也解释了为什么课程会横跨 imitation learning、model-based、RL for LLMs、robotics 等多个主题，因为这些方向分别在解决不同瓶颈。

\subsection{本章小结}

学习深度强化学习有四个层面的理由：它能处理监督学习难以处理的决策后果问题；它支撑大量现代 AI 系统；它提供了学习于经验的形式化框架；同时它仍然有大量基础与应用研究问题值得解决。

\section{如何把行为写成数据：状态、观测、动作、轨迹与奖励}

\subsection{顺序决策问题的最小词汇表}

讲者接着把 RL 的基础记号完整展开。要把一个问题变成强化学习问题，至少要写清这几个对象：
\begin{itemize}
\item state \(s_t\)：时刻 \(t\) 世界的状态；
\item observation \(o_t\)：智能体在时刻 \(t\) 实际看到的信息；
\item action \(a_t\)：在时刻 \(t\) 做出的决策；
\item trajectory \(\tau\)：状态 / 观测与动作构成的时序序列；
\item reward \(r(s,a)\)：状态动作对有多好。
\end{itemize}

\begin{figure}[H]
\centering
\includegraphics[width=0.95\textwidth]{pdf_pages/page-32.png}
\caption{把经验表示为数据的基本对象：状态、观测、动作、轨迹与奖励。}
\end{figure}

其中最容易混淆的是 state 和 observation。讲者专门提醒，两者并不总相同。状态是“足够描述未来演化”的世界变量；观测则只是智能体拿到的输入，有时它不完整，只是对状态的部分投影。

\subsection{状态与观测：Markov 性质为何重要}

当下一个状态只依赖当前状态、当前动作以及随机性，而不依赖更久远的历史时，我们说系统满足 Markov property。这个性质之所以重要，是因为它把复杂的历史依赖压缩进当前状态中，让决策问题变得可递推、可估计。

\begin{figure}[H]
\centering
\includegraphics[width=0.92\textwidth]{pdf_pages/page-33.png}
\caption{状态、观测与系统动力学的关系：Markov 性质要求未来由当前状态和动作决定。}
\end{figure}

但在很多现实问题中，智能体拿不到完整状态，只能看到观测。例如对话系统只能看到最近对话历史，机器人可能只能拿到图像与传感器读数。这时如果只把 \(o_t\) 喂给策略，信息可能不够，往往要给策略加记忆，让它基于一段观测历史作决策。

\subsection{轨迹与奖励：从单步判断走向长期效果}

轨迹可以写成
\[
\tau=(s_1,a_1,s_2,a_2,\ldots,s_T,a_T).
\]

符号说明如下：
\begin{itemize}
\item \(T\)：轨迹长度，可以固定也可以变化；
\item \(s_t,a_t\)：第 \(t\) 步的状态与动作；
\item 整条轨迹是一次 policy roll-out，也常被叫作一个 episode。
\end{itemize}

奖励函数则用来把“想要什么行为”写成数值目标。例如：
\begin{itemize}
\item 毛巾挂到钩子上给高分；
\item 用户给出 upvote 给高分；
\item 机器人向前移动得更快给高分。
\end{itemize}

这一步是强化学习建模最关键也最脆弱的地方。因为奖励写得不好，系统就会学到偏离真实意图但在数值上看起来更优的行为。

\subsection{策略：怎样用神经网络表示行为}

策略 \(\pi_\theta\) 负责把状态或观测映射到动作分布。第一讲里先不给出具体网络结构，而是强调策略是“行为的参数化表示”。这让后续所有算法都可以统一理解为：调整参数 \(\theta\)，让策略生成更好的轨迹。

讲者还特别提出 stochastic policies 的动机。引入随机性并不是装饰，而是因为：
\begin{itemize}
\item 需要探索时，系统必须尝试不同动作；
\item 真实示范数据和现实世界本身都常常是多峰的、随机的；
\item 用分布而非单点，更能表示行为不确定性。
\end{itemize}

\subsection{本章小结}

把问题写成 RL，需要先定义数据对象，再定义行为对象。状态与观测区分了“世界真实变量”和“智能体实际可见信息”；轨迹与奖励把问题从单步预测升级为长期结果优化；策略则把行为变成可学习的参数化对象。

\section{强化学习的目标函数与算法版图}

\subsection{长期目标：最大化期望累计奖励}

讲者把 RL 的核心目标写成：
\[
\max_\theta \; \mathbb{E}_{\tau\sim p_\theta(\tau)}\left[\sum_{t=1}^{T} r(s_t,a_t)\right].
\]

这里每个符号都很重要：
\begin{itemize}
\item \(\theta\)：策略参数；
\item \(p_\theta(\tau)\)：由当前策略与环境动力学共同诱导的轨迹分布；
\item \(\sum_{t=1}^{T} r(s_t,a_t)\)：一条轨迹上的累计回报；
\item \(\mathbb{E}\)：我们优化的是在轨迹分布下的平均表现，而不是单次 rollout 的运气。
\end{itemize}

这条公式把前面所有概念串了起来。策略不是为了某一步动作像不像专家，而是为了让长期结果更好。

\subsection{为什么常常需要随机策略}

第一讲给出的随机策略理由很朴素，但非常根本：
\begin{itemize}
\item 没有随机性，就很难主动探索没见过的动作；
\item 如果现实数据来自多个示范者或多种成功方式，确定性输出容易退化成“平均动作”；
\item 后续的很多梯度型算法，本来就是基于动作分布来推导的。
\end{itemize}

因此“策略是分布”并不是技术细节，而是深度强化学习方法设计的基础假设之一。

\subsection{算法类型：课程接下来会走哪几条线}

在课程结尾，讲者给出一个非常有用的算法地图。不同算法都服务于同一个目标函数，但会在监督来源、数据来源、是否显式估计价值、是否学习动力学模型等方面做不同权衡。

\begin{figure}[H]
\centering
\includegraphics[width=0.93\textwidth]{pdf_pages/page-41.png}
\caption{强化学习算法版图：模仿学习、策略梯度、actor-critic、value-based 与 model-based。}
\end{figure}

这几类方法可以先做粗略理解：
\begin{itemize}
\item imitation learning：直接模仿高回报行为；
\item policy gradients：直接对上面的期望目标求梯度；
\item actor-critic：一边学策略，一边学价值评估来帮助改进策略；
\item value-based：先学“什么动作值钱”，再由价值导出策略；
\item model-based：显式学习环境动力学，再做规划或利用模型辅助学习。
\end{itemize}

讲者没有说某条路线绝对更好，而是强调方法选择取决于数据是否便宜、奖励是否清晰、环境是否可控、是否需要高数据效率等条件。

\begin{warningbox}{为什么会有这么多 RL 算法}
因为强化学习的难点不是单一的。有人缺数据，有人缺稳定性，有人缺可探索性，有人环境危险无法在线试错，所以算法会围绕不同瓶颈做取舍。
\end{warningbox}

\subsection{本章小结}

强化学习的统一目标是最大化期望累计奖励，但实现这件事可以走多条路线。第一讲的作用，就是先把整个方法空间摆出来，后续每一讲再把其中一类方法拆开讲透。

\section{总结与延伸}

Lecture 1 完成了两件事。第一，它把 CS224R 这门课的学习框架交代清楚了：课程强调实现、实验和项目，不打算停留在高层概念综述。第二，它用最少但最关键的记号搭建了 RL 的语言体系：状态、观测、动作、轨迹、奖励、策略、期望累计回报。

如果用一句话概括本讲的核心，那么就是：\textbf{强化学习研究的是会影响未来数据分布的行为优化问题。} 这使它天然不同于静态监督学习，也解释了为什么后续会出现 imitation learning、policy gradient、actor-critic、value-based 和 model-based 这些不同路线。

从学习路径看，第一讲最值得反复消化的不是 logistics，而是三个基本判断：
\begin{itemize}
\item 行为学习问题和静态预测问题在数据生成机制上不同；
\item 建模前必须先分清 state、observation、reward 和 trajectory；
\item 所有后续算法，本质上都在尝试更有效地优化长期回报目标。
\end{itemize}

后续课程自然会沿着这个框架推进：先从 imitation learning 这样“监督信号更强”的方法出发，再进入真正依赖在线试错的 reinforcement learning 算法。

\end{document}
